% chapters/ch09.tex —— 第 9 章
% 体例：具体失败 → 设计问题 → 机制解释 → 对照实验 → 条件变化 → 选择原则
% 数字身份纪律：AFlow 与 ADAS 的数字一律 \cited{}（出处见 source-map 第 9 章各条）；
% 由引用数字算出的量用 \calcnum{}。E06 尚未运行，本章不报告它的任何数字。
\chapter{工作流与多 Agent 结构的演化}
\label{ch09}

\chaptergoal{应该改一个节点，还是改组织这些节点的方式？}
\chapterprops{本章沿信息流解释结构改动，分别计算搜索投入与部署资源，再用相同总预算的多采样对照检查结构增量。E06保留为尚未运行的B级受限设计。}
\chapterfig{图~\ref{fig:ch09-workflow}：三种组织方式及各自的资源账}
\chapterpitfall{把结构变复杂当成找到了更好的结构；把搜索成本排除在比较之外；
把跨模型能用当成跨模型无损}

% ---------- 1. 具体失败 ----------
\begin{failurebox}
AFlow 用 Claude-3.5-sonnet 作为优化器，以 GPT-4o-mini 的执行反馈搜索出的工作流，换用 DeepSeek-V2.5 执行时报告
\cited{0.9390} 的 HumanEval 分数，而 GPT-4o 直接调用是 \cited{0.9389}——
论文报告前者的测试执行费用为后者的 \cited{4.55\%}\cite{aflow-2410}。

两个分数的点估计接近。统计等效需要相应比较证据，费用比也有明确范围。
那个 \cited{4.55\%} 只统计\emph{测试执行 token 按当时模型价格折算的美元费用}。
该实验的 MCTS 最大轮数设为 \cited{20}、每轮把候选工作流在验证集上跑 \cited{5} 次\cite{aflow-2410}，
搜索费用没有进入这张测试执行费用表；未报告部署复用量时，不能说已完成摊销。

同一次比较里还有第二个缺口：
作者列出的基线里已经有多尝试方法（CoT 自一致性 \cited{5} 个答案、
MedPrompt \cited{3} 个答案 \cited{5} 票、Self-Refine、MultiPersona），
但``给一个扁平采样器与结构搜索同等预算、取 best-of-N''这一组对照并未出现\cite{aflow-2410}。
也就是说：\textbf{我们不知道这 \cited{4.55\%} 换来的分数，
是不是同样多的钱用来多采样几次也能拿到。}
\end{failurebox}

\designquestion{判据不是``搜出来的工作流分数更高''，而是：
\textbf{在同等总预算下，把钱花在搜索结构上，是否比花在同一次推理的多采样上更值}。
回答它需要同时保存搜索投入、部署资源和任务结果。}

第7章以指令、示例和上下文为主要对象，第8章讨论可复用的可执行片段。
本章将\emph{节点之间的连接}作为显式搜索对象；程序优化也可能跨越这些边界。
这一层同时允许提示与连接变化，可能增加候选组合和评价负担，
因此也最容易在预算口径上出问题。

% ---------- 2. 机制解释 ----------
\section{工作流的表示：节点、边与算子}
\label{sec:ch09-repr}

AFlow 把一个工作流写成节点序列 $N$ 与用代码表示的边 $E$，
每个节点是一个带提示、温度与格式的模型调用，边决定调用顺序、条件与循环\cite{aflow-2410}。
优化的目标是

\begin{equation}
W^{*}=\arg\max_{W\in\mathcal S} G(W,T),
\label{eq:ch09-aflow}
\end{equation}

其中 $G$ 是工作流 $W$ 在任务 $T$ 上的得分\cite{aflow-2410}。

用代码表示边带来一个直接的后果：顺序、分支、循环、并行与路由
都能被同一个表示装下，不需要为每种结构单独设计一套语法。
ADAS 把这一点推到了头——既然多数编程语言是图灵完备的，
``用代码定义 agent''在理论上能表达任何 agentic system，
包括新的提示、工具用法、工作流及其组合\cite{adas-2408}。

\begin{principlebox}
\textbf{表达能力强不等于搜索得动。}
通用代码表示具有很强的表达能力，但实际工具接口、运行预算与安全限制仍会约束可执行集合，搜索器也仍有先验。
AFlow 因此把搜索交给 MCTS：树的一个节点是\emph{一整个工作流}，
而不是工作流里的一个节点\cite{aflow-2410}。
评估单位随之变成``把整个工作流在验证集上跑一遍''——
提示优化也可能评估整个多模块程序，因此不能预定这里贵一个数量级，费用需要按实际任务和候选结算。
\end{principlebox}

ADAS 给出另一个部件：一个持续增长的 archive。
被发现的 agent 连同它的指标一起入档，元 agent 参考这份档案设计下一个\cite{adas-2408}。
这就是结构搜索里的``经验''——注意它存的是\emph{被评估过的候选}，
是服务于搜索过程的持久状态；它与部署时读取的经验用途不同，是否跨搜索复用取决于设计，不能由名称判定它不持久。

\section{五种结构原语各自承担什么}
\label{sec:ch09-primitives}

用代码表示边的好处是表达能力强，代价是``结构''这个词变得太大。
把它拆开看，本章用五种常见控制原语拆解工作流，不声称穷尽全部结构；它们的功能不同，
也因此带来不同的失败方式。

\begin{table}[htbp]
\centering
\caption{五种结构原语：各自承担的功能、新增的失败方式与必须记录的量}
\label{tab:ch09-primitives}
\begin{tabular}{>{\raggedright\arraybackslash}p{1.6cm}>{\raggedright\arraybackslash}p{3.4cm}>{\raggedright\arraybackslash}p{3.5cm}>{\raggedright\arraybackslash}p{4.0cm}}
\toprule
原语 & 承担的功能 & 新增的失败方式 & 必须记录的量 \\
\midrule
顺序 & 把一个结果喂给下一个节点 & 上游的错误被下游当作事实继续推导 & 每一跳的输入输出对 \\
分支 & 按条件选择不同路径 & 条件错误导致错误路径；无日志时难追查 & 每次判定的条件与选中分支 \\
循环 & 在达到条件前反复修改 & 没有可判定出口时调用数无上界 & 迭代次数与每轮的改动量 \\
并行 & 同一问题取多条可能相关的尝试 & 新增生成与聚合费用；相关错误可能保留 & 路数、每路结果、聚合规则 \\
路由 & 按问题类型分发到不同子流程 & 分发错时错误被归因到被选中的子流程 & 分发判据与命中分布 \\
\bottomrule
\end{tabular}
\end{table}

\begin{sloppypar}
表~\ref{tab:ch09-primitives} 的第二列说明一件容易被说反的事：
\textbf{这些原语改变信息流与计算分配，是否提高任务能力需要对照。}
顺序把不确定性留在下游，分支把它变成一次判定，循环把它换成更多调用，
并行把它摊到多条路上，路由把它交给一次分类。
因此``加了某个结构所以分数上升''这句话缺少一个中间环节：
必须说明是哪一个原语把哪一种不确定性换掉了，以及换来的代价是什么。
\end{sloppypar}

表~\ref{tab:ch09-primitives} 的后两列把失败方式与记录位置对应起来。例如，分支条件选错后，最后答案的零分无法说明问题出在路由还是求解；保存分支判定及下游输入，才能沿调用链定位。质量、执行路径与资源需要一起读。

\subsection{合法结构的约束：搜索空间必须删掉什么}
\label{sec:ch09-legal}

图灵完备的表示能装下任何结构，包括运行不起来的那些。
一个结构搜索器如果不先删掉一部分结构，就会把预算花在证明它们不行上。
下面四条约束分别处理终止、状态、评价与接口，避免无界执行或无法解释的结果：

\begin{enumerate}
  \item \textbf{循环必须有可判定的出口与强制上界。}
        验收结果可触发提前退出，计数器或运行预算还须保证始终未通过时也会停止。
        允许``改到满意为止''等于允许调用数无上界——
        第 \ref{ch02} 章的 episode 预算在这里变成结构层面的预算。
  \item \textbf{并行共享状态必须有明确同步语义。}
        缺少同步时可能出现调度依赖；可以使用隔离副本、事务或明确的合并规则。随机性本身不使配对比较失效，但状态耦合须进入设计与分析。
  \item \textbf{验收器不在结构的可写侧。}
        这是第 \ref{ch01} 章公理一（验证器冻结）在结构层的同一条约束：
        本章固定共同判据；若另行研究评价器更新，旧分差不能直接解释为同尺度改进，应保留共同外部参照。
  \item \textbf{每个节点的输出必须有类型契约。}
        节点可以传文本或结构化对象；解析器能捕获部分形状错误，语义错误仍可能传到下游。
        契约缺失时，失败会被归因到最后一个节点，而真正的起因在上游。
\end{enumerate}

\begin{principlebox}
\textbf{表达能力、合法性检查与搜索效率要分别说明。}
ADAS 与 AFlow 都用代码表示候选。AFlow 提供算子来帮助搜索，但也允许通过 Custom 构造节点，并报告不使用预定义算子的消融；不能把它描述为只允许七种固定结构。预定义接口可以减少部分无效候选，却不保证整个程序一定终止或正确。反过来，通用表示也可以配合运行时预算和校验器。
\end{principlebox}

\subsection{提示优化与结构优化：两个不同的搜索空间}
\label{sec:ch09-space}

\begin{table}[htbp]
\centering
\caption{两类优化的搜索空间与评估代价}
\label{tab:ch09-space}
\begin{tabular}{>{\raggedright\arraybackslash}p{2.8cm}>{\raggedright\arraybackslash}p{5.0cm}>{\raggedright\arraybackslash}p{5.1cm}}
\toprule
 & 提示 / 示例优化（第 7 章） & 结构优化（本章） \\
\midrule
搜索对象 & 模块指令、示例与上下文 & 节点之间的边与算子组合 \\
候选评估 & 可能评估完整多模块程序 & 按完整候选及其路径评估 \\
行为影响 & 可影响本模块及其下游输入 & 可改变可达路径、调用顺序与资源分配 \\
失败形态 & 塌缩、记住题目 & 调用数膨胀、等待时间变长 \\
\bottomrule
\end{tabular}
\end{table}

表~\ref{tab:ch09-space} 最后一行的差别最容易被忽略：
两类优化都可能改变分数、调用数和等待时间。提示可以诱导额外工具调用，结构可以改变答案质量，不能将两种失败形态绝对分开。
开头的费用案例需要按同样的原则分解。

AFlow 的算子集合给出结构搜索的原子：
Generate、Format、Review \& Revise、Ensemble、Test、Programmer、Custom\cite{aflow-2410}。
算子可以封装多个调用或条件执行，增加一个算子不等于固定增加一次模型调用。

\section{候选工作流的生成与局部修改}
\label{sec:ch09-search}

AFlow 的一轮迭代是四步：Select（软混合概率选择）、
Expand（由模型优化器改提示与代码边生成子工作流）、
Evaluate（验证集上测 \cited{5} 次）、
Backpropagate（记录性能与这一改的成败）；
终止条件是最优若干候选连续若干轮无改进，或总轮数到 \cited{20}\cite{aflow-2410}。

\begin{sloppypar}
\textbf{先保留原文中的一个参数歧义。}
正文式(3)说明采用 $\lambda=\cited{0.2}$、$\alpha=\cited{0.4}$，但附录 A.6 详细算法第36行写作 $\lambda=\cited{0.4}$、$\alpha=\cited{0.2}$\cite{aflow-2410}。前者控制均匀探索权重，后者控制分数对选择概率的影响，交换它们不是无关紧要的排版变化。固定公开实现又采用另一组默认值，见下一小节。论文说明、代码默认项和历史运行有效配置需要分别记录，原表参数尚缺检查点或运行配置的直接绑定。

\textbf{Backpropagate 这一步值得单独看。}
它把候选相对父候选的表现及修改记录写回树里，形成关于修改效果的反馈。
它与第 6 章的准入判决、第 8 章的入库判决是同一类东西：
都是``这次更新值不值得''的显式记录。
这份经验主要用于搜索选择；原文没有据此规定搜索结束必须销毁，跨搜索复用需要额外设计和验证。
\end{sloppypar}

\subsection{分数变换也属于选择参数}
\label{sec:ch09-score-scale}

固定公开提交的选择函数默认 $\alpha=\cited{0.2}$、$\lambda=\cited{0.3}$，先将读入的得分乘以\cited{100}，再计算指数权重并与均匀分布混合\cite{aflow-code-2025}。对非空且得分有限的候选集，令 $0\le\lambda\le1$，用原始得分 $s_i$ 表示这条路径，其概率为
\begin{equation}
 p_i=\frac{\lambda}{n}+(1-\lambda)
 \frac{\exp(\alpha\widetilde s_i)}
      {\sum_{j=1}^{n}\exp(\alpha\widetilde s_j)},
 \qquad \widetilde s_i=100s_i.
 \label{eq:ch09-reference-selection}
\end{equation}
实现先减去最大分数再取指数，得到相同概率并改善数值稳定性。均匀项只在已经进入当前父候选集合的工作流之间分配，不表示在整个代码搜索空间中均匀抽样。

以本书构造的两个原始得分\calcnum{0.80}、\calcnum{0.81}为例，代入这份参考实现的默认值，较高分者被抽到的概率约为\calcnum{0.5349}；若遗漏乘以百的步骤，就会算成约\calcnum{0.5003}。这两个数是公式计算，不是一次搜索实测。系数只有连同分数单位和变换一起保存，才能说明选择偏好。

这份仓库说明包含从 MetaGPT 迁移后的实现边界\cite{aflow-code-2025}。上述默认值能够定位公开脚本的行为，尚不能替代原论文两组参数所对应的历史运行配置。复现记录应保留原始得分、变换、候选集合与最终抽样概率。

\subsection{搜索看过哪些题，决定它学到了什么}
\label{sec:ch09-validationselection}

AFlow 的初始化还有一个与预算和泛化同时相关的细节。原文说明，先用固定随机种子\cited{42}把数据划分为\cited{20\%}验证集与\cited{80\%}测试集；随后在验证集上将空白模板运行\cited{5}次，选出得分方差较高的题目作为最终验证子集\cite{aflow-2410}。这一步发生在工作流迭代之前，并且使用模型反馈决定后续搜索主要看哪些题。

高方差题可能更有区分力：对于始终答对或始终答错的题，一次小结构修改可能很难在有限重复中显示差别；处在边界的题则可能更敏感。但这也改变了搜索所见任务的分布。最终验证均值代表的是筛选后的子集，不能不加说明地当成原验证分布的无偏估计。独立测试仍可评价选中结构在预定测试范围内的表现，前提是测试未再被拿来反复挑选结构。

这项筛选不是因为使用了反馈就必然构成泄漏。关键在于反馈属于哪一个分区、在哪个决策中使用，以及最终主张针对哪个范围。训练或验证数据用于选择候选，本来就是搜索的一部分；把测试结果也拿来选赢家后仍称其为未参与选择的最终检验，才混淆了角色。ADAS 按来源测试表现挑选转移候选的例子，正说明同一个名字为“测试”的集合可能在流程中承担了选择作用\cite{adas-2408}。

初始化的五次运行也要计入搜索成本。若最终报告只从第一轮候选生成开始记账，选题本身就被漏掉。类似地，为设定停止轮数、比较提示模板或调节选择概率而做的探索，都可能消耗资源；它们是否进入当前方法的完整成本，应按预先定义的研究与部署范围说明。没有必要把无法恢复的历史探索硬凑成精确数字，但应保留这部分为未计费用，而不是宣称全部搜索成本已知。

自适应选题还会影响候选之间的可比性。若每个候选都在不同的“最有区分力”题集上评价，均值差异同时包含了结构变化和题目变化。可以固定一份共同评价子集，也可以设计有明确估计方法的自适应分配，但不能简单比较不同题集的原始均分后宣布后者更优。评价记录至少需要保留逐题标识，以便事后发现覆盖差异。

近期预印本中的 EvalAlloc 将搜索反馈与候选确认分成两个阶段\cite{evalalloc-2610}。它先在 SEARCH 题集上完整评价预热候选，建立共享的高信息量题目核心与共同评分参照。随后用稀疏观察为进化搜索提供反馈，搜索期间保持题目参数与评分参照固定。搜索结束后，对有效候选逐轮扩大共享的评价题目前缀并淘汰候选，复用已有逐题结果，只补跑缺失项。进入最后一轮的候选补齐全部 SEARCH 观测后，再按实际结果选出最终程序。这套协议中的 SEARCH 同时用于进化与最终候选选择，DEV 仅作审计，TEST 只评价选出的程序。其预算分配结论仍受所用框架、任务与评价单元约束。

重复评价的作用同样有限。同一份验证集上运行五次，可以观察采样波动，却不能凭此创造五份独立任务分布。若某个候选在这份固定集合上具有特有优势，多次重复仍会保留这种优势。报告时应区分任务抽样的不确定性、模型采样的不确定性和搜索选择的不确定性；只对最终候选的五个运行均值求区间，无法自动覆盖整个自适应搜索过程的选择影响。

档案可以帮助保留这些条件。一个可复用的搜索条目不应只有“结构甲得分最高”，还应写明它在哪个验证版本、哪一批题、什么后端和预算下被选中。再次搜索时，档案中的失败可能提示接口风险，也可能只是旧题分布的局部结果。让新搜索者读取带范围的记录，才有机会使用历史信息而不把旧偏差不断放大。

\section{多 Agent 的角色、通信与协作结构}
\label{sec:ch09-multi}

上面的算子换个说法就是角色：Review \& Revise 是一个评审者，
Ensemble 聚合多个候选，候选生成是否并行要看执行代码，Programmer 是一个会写代码的节点。
在本章采用的表示中，可以把多 Agent 协作展开为工作流：
角色对应节点，通信对应边。这是分析信息流的方法，不要求所有多 Agent 系统都采用预先固定的角色图。

\begin{table}[htbp]
\centering
\caption{AFlow 在六个基准上的整体结果（GPT-4o-mini 执行，原论文 Table~1）}
\label{tab:ch09-aflow}
\begin{tabular}{>{\raggedright\arraybackslash}p{2.4cm}cccccc}
\toprule
 & HotpotQA & DROP & HumanEval & MBPP & GSM8K & MATH \\
\midrule
AFlow & \cited{73.5} & \cited{80.6} & \cited{94.7} & \cited{83.4} & \cited{93.5} & \cited{56.2} \\
\bottomrule
\end{tabular}

\vspace{2mm}
\begin{footnotesize}
数字身份为[引用]，来自原论文 Table~1\cite{aflow-2410}。每种方法测试 \cited{3} 次后取均值。
前两项为 F1，中间两项为 pass@1，后两项为解题率，数值沿用原表的显示尺度。六基准等权平均为 \cited{80.3}。
原文报告相对最强手工基线与 ADAS 的相对提升约为 \cited{5.7\%} 与 \cited{19.5\%}，不是百分点；按表中均值相减分别为 \calcnum{4.3} 与 \calcnum{13.1} 个分数点。
这些数值保留原论文的报告身份。比较工作流的资源效率，还需并列报告调用数与等待时间。
\end{footnotesize}
\end{table}

表~\ref{tab:ch09-aflow}提供了整体质量比较。
它能支持的说法是``在这六个基准上，搜出来的工作流分数更高''；
它不能支持``这些分数是用更少的资源拿到的''——
因为那张表没有资源这一栏。

% figures/ch09-workflow.tex —— 第 9 章核心图（全书引用入口）
% 用法：在 book/chapters/ch09.tex 中 % figures/ch09-workflow.tex —— 第 9 章核心图（全书引用入口）
% 用法：在 book/chapters/ch09.tex 中 % figures/ch09-workflow.tex —— 第 9 章核心图（全书引用入口）
% 用法：在 book/chapters/ch09.tex 中 \input{figures/ch09-workflow}
% 图主体在 ch09-workflow-body.tex，便于单独编译与复用。
\begin{figure}[htbp]
\centering
\input{figures/ch09-workflow-body}
\caption{顺序、带评审与并行集成的控制流示意。图中节点和分支不同，具体调用量、时延与质量需按实现测量；并行不保证等待时间不变，评审也不保证纠错。即使最终分数相同，执行过程与资源账仍可能不同。}
\label{fig:ch09-workflow}
\end{figure}

% 图主体在 ch09-workflow-body.tex，便于单独编译与复用。
\begin{figure}[htbp]
\centering
% Three full-size panels; no transform-shape reduction.
\begin{tikzpicture}[
 nd/.style={draw,rounded corners=2pt,text width=18mm,minimum height=8mm,align=center,font=\small},
 io/.style={nd,fill=gray!10},
 arrow/.style={-{Stealth[length=1.8mm]},thick},
 note/.style={font=\small,align=left,text width=43mm}
]
\node[font=\bfseries\small] at (0,1) {顺序};
\node[io] (s1) at (0,0) {生成};
\node[nd] (s2) at (0,-1.2) {求解};
\node[io] (s3) at (0,-2.4) {输出};
\draw[arrow] (s1)--(s2);\draw[arrow] (s2)--(s3);
\node[note,anchor=north] at (0,-3.2) {按实际调用与\\执行事件记账。\\等待时间沿串行路径累加。};
\node[font=\bfseries\small] at (5,1) {带评审};
\node[io] (r1) at (5,0) {生成与求解};
\node[nd,fill=orange!12] (r2) at (5,-1.2) {评审};
\node[nd] (r3) at (6.1,-2.4) {修正};
\node[io] (r4) at (5,-3.6) {输出};
\draw[arrow] (r1)--(r2);\draw[arrow] (r2)--(r3);
\draw[arrow] (r3)--(r4);\draw[arrow] (r2.west)--++(-1.2,0)|-(r4.west);
\node[note,anchor=north] at (5,-4.4) {评审与修正分别计费。\\时间取决于实际返回分支。};
\node[font=\bfseries\small] at (10,1) {并行集成};
\node[io] (p0) at (10,0) {输入};
\node[nd,fill=blue!9] (p1) at (8.8,-1.5) {求解甲};
\node[nd,fill=blue!9] (p2) at (11.2,-1.5) {求解乙};
\node[io] (p3) at (10,-3) {聚合};
\draw[arrow] (p0)--(p1);\draw[arrow] (p0)--(p2);
\draw[arrow] (p1)--(p3);\draw[arrow] (p2)--(p3);
\node[note,anchor=north] at (10,-3.8) {各路与聚合均计费。\\等待受并发和排队影响。\\质量取决于错误相关性。};
\end{tikzpicture}

\caption{顺序、带评审与并行集成的控制流示意。图中节点和分支不同，具体调用量、时延与质量需按实现测量；并行不保证等待时间不变，评审也不保证纠错。即使最终分数相同，执行过程与资源账仍可能不同。}
\label{fig:ch09-workflow}
\end{figure}

% 图主体在 ch09-workflow-body.tex，便于单独编译与复用。
\begin{figure}[htbp]
\centering
% Three full-size panels; no transform-shape reduction.
\begin{tikzpicture}[
 nd/.style={draw,rounded corners=2pt,text width=18mm,minimum height=8mm,align=center,font=\small},
 io/.style={nd,fill=gray!10},
 arrow/.style={-{Stealth[length=1.8mm]},thick},
 note/.style={font=\small,align=left,text width=43mm}
]
\node[font=\bfseries\small] at (0,1) {顺序};
\node[io] (s1) at (0,0) {生成};
\node[nd] (s2) at (0,-1.2) {求解};
\node[io] (s3) at (0,-2.4) {输出};
\draw[arrow] (s1)--(s2);\draw[arrow] (s2)--(s3);
\node[note,anchor=north] at (0,-3.2) {按实际调用与\\执行事件记账。\\等待时间沿串行路径累加。};
\node[font=\bfseries\small] at (5,1) {带评审};
\node[io] (r1) at (5,0) {生成与求解};
\node[nd,fill=orange!12] (r2) at (5,-1.2) {评审};
\node[nd] (r3) at (6.1,-2.4) {修正};
\node[io] (r4) at (5,-3.6) {输出};
\draw[arrow] (r1)--(r2);\draw[arrow] (r2)--(r3);
\draw[arrow] (r3)--(r4);\draw[arrow] (r2.west)--++(-1.2,0)|-(r4.west);
\node[note,anchor=north] at (5,-4.4) {评审与修正分别计费。\\时间取决于实际返回分支。};
\node[font=\bfseries\small] at (10,1) {并行集成};
\node[io] (p0) at (10,0) {输入};
\node[nd,fill=blue!9] (p1) at (8.8,-1.5) {求解甲};
\node[nd,fill=blue!9] (p2) at (11.2,-1.5) {求解乙};
\node[io] (p3) at (10,-3) {聚合};
\draw[arrow] (p0)--(p1);\draw[arrow] (p0)--(p2);
\draw[arrow] (p1)--(p3);\draw[arrow] (p2)--(p3);
\node[note,anchor=north] at (10,-3.8) {各路与聚合均计费。\\等待受并发和排队影响。\\质量取决于错误相关性。};
\end{tikzpicture}

\caption{顺序、带评审与并行集成的控制流示意。图中节点和分支不同，具体调用量、时延与质量需按实现测量；并行不保证等待时间不变，评审也不保证纠错。即使最终分数相同，执行过程与资源账仍可能不同。}
\label{fig:ch09-workflow}
\end{figure}


图~\ref{fig:ch09-workflow} 把缺的那一栏画出来：
三种组织方式即使拿到同一个分数，
它们的调用数与等待时间仍可能不同。
只看分数时，三者不可区分。


\subsection{多角色的增量来自信息，而不来自称呼}
\label{sec:ch09-roleinformation}

给同一个模型分别设置“专家”“反对者”“裁判”三个角色，可能改变它生成的内容，但角色名本身不提供新的观测。若所有角色都读到同一份错误检索材料，又不能访问其他来源，它们很可能围绕错误前提进行精致讨论。这里不需要假定它们总是相关，也不需要否定角色提示的作用；只需要把声称的机制说清楚：增加的是表达视角、独立采样、额外工具，还是有来源的外部证据。

一个反对者如果能够运行独立测试，就多了一条与求解者解释不同的检查通道；如果只能要求“再仔细想想”，则主要增加一次模型推理。二者都可能有用，成本与失败方式却不同。前者可能受测试覆盖不足限制，后者可能沿用同样的错误前提。评价时可以固定角色提示只改变信息来源，也可以固定信息来源只改变分工，从而分别考察哪一项带来可重复的增量。

通信内容也不宜只按长度理解。两个节点各收到一千字，信息可能完全不同：一份包含原始测试输出，另一份只有上游的结论摘要。摘要可以降低成本，却可能删除异常条件、冲突证据或版本标识。若使用摘要结构与直接共享原始轨迹的结构比较，应同时保存原文与摘要，说明哪些字段被保留，而非只报告 token 节省。信息损失未必有害，关键是它是否改变了任务需要的判断依据。

对抗式分工还有一个常见误区：让批评者必须找出错误。这样的提示可能提高问题发现率，也可能鼓励它把正确答案改坏。比较应同时观察错误候选被纠正、正确候选被破坏两个方向。在只保存“批评成功案例”的档案中，后一个方向很容易消失。若任务对误改代价很高，结构可能需要把建议与实际修改分开，再要求一个可执行条件支持修改。

协作结构的消融则需要考虑接口重连。删掉一个评审节点后，直接把它的输入传给下游，可能改变格式；如果删掉的节点同时完成格式转换，分数下降未必来自失去评审能力。较清楚的设计是保留必要的确定性格式转换，只移除要研究的判断功能。反过来，若研究的问题就是整个复合节点是否值得保留，那么连同转换一起删除也合理，但应按复合干预解释结果。

有些结构只在少数任务上被触发。比如困难路由仅服务罕见长题，整体均分上的作用可能很小，却能改变那一类任务的完成率。报告全量结果的同时，可以给预先定义的类别结果；不能只在看到收益后挑一个有利子组。实际触发率、路由失败率和后备路径成本一起记录，才能解释某个角色为什么在多数题上看不见，却仍然有部署价值。

还要注意部件之间的相互作用。假设检查者输出一份精确错误日志，修订者才会据此有效修改；单独增加检查者而不使用日志，最终答案可能不变，单独增加修订者而不给日志，也可能无效。两者合用却可能产生增量。这里的四种配置分别是都无、仅检查、仅修订、检查加修订，能够帮助识别“合用效果是否不同于各自效果相加”。这是一种比较设计，不是本书已观测到的模型结果。

若以四种配置的质量均值记为 $Q_{00},Q_{10},Q_{01},Q_{11}$，可考察交互差 $I=Q_{11}-Q_{10}-Q_{01}+Q_{00}$。非零点值本身仍需不确定性分析；即使得到可靠差异，也只描述当前尺度上的交互，不能直接推断内部认知机制。各配置的资源也不同，应同时记录，避免把额外调用带来的收益全部归为部件协同。

反过来，从完整结构里分别删除一个部件，只回答该部件在其他部件仍然存在时的条件贡献。它没有覆盖“两个部件一起删除后是否可以换一种简单策略”的问题。若工作流有强互补或冗余，逐项消融就不足以给每个节点分配一个唯一功劳值。可以依据机制挑选少量联合消融，而不必穷举所有组合；选择依据和未覆盖范围需要写清楚。

失败路径也可能需要联合检查。正常答案不触发后备路由时，删除后备节点看不到质量变化；这不能证明节点无用。可在隔离环境下构造服务超时或格式错误，检查后备行为是否符合协议。这种故障注入测量的是工程恢复，不应与自然任务的成功率混在同一分母，也不能将人为制造的高故障率当成真实服务故障率。

最终，是否增加一个角色可以落实为很具体的问题：它得到什么此前没有的信息，它能作什么此前不能作的检查，它的输出会改变哪个决定，删除它之后哪些条件保持相同。若这些问题答不清，继续增加角色很可能只增加解释上的复杂度。若答得清，即使结构很简单，也已经具备了被验证和被改进的起点。

\section{从 AFlow 示例还原工作流}
\label{sec:ch09-codecase}

AFlow 的 GSM8K 案例适合练习如何读工作流，因为它把候选代码与搜索路径放在了相邻图中。原论文 Figure~5(B) 展示的程序先生成 \cited{5} 个候选答案，再执行集成，最后调用 Programmer 检查集成结果；如果检查返回的输出非空，就采用该输出，否则保留集成结果\cite{aflow-2410}。这段程序包含生成、汇总、检查和后备返回，读者很容易将它脑补成“多路并行求解后由验证器挑出正确答案”。原代码支持的事实比这句话窄。

下面是依据v4 Figure~5(B)独立改写的最小伪代码。它保留调用顺序与返回条件，省略类定义、提示常量及成本对象；不是原实现摘录，也不是可直接复现实验的代码。

\noindent\begin{minipage}{\linewidth}
\begin{lstlisting}[language=Python,caption={AFlow GSM8K示例的控制流教学改写},label={lst:ch09-aflow},firstnumber=1]
async def solve(problem):
    answers = []
    for attempt in range(5):
        answer = await generate(problem)
        answers.append(answer)
    combined = await ensemble(problem, answers)
    checked_output = await programmer(problem, combined)
    if checked_output:
        return checked_output
    return combined
\end{lstlisting}
\end{minipage}

第3--5行逐个等待五次生成，第6行等待集成，第7行才开始Programmer检查。若第8行条件为真，第9行返回检查输出；否则落到第10行。顶层最多显示七次算子入口，底层请求数仍需打开各算子的实现核对。第8行仅测试输出非空，没有读取标准答案，也没有证明正确。

首先，生成循环在每一次迭代中等待当前调用结束，然后才进入下一次。函数采用异步语法，不等于循环里的请求已经同时发出。若外部服务内部另有调度，那是另一个执行层；图中这段调用方代码本身没有把多个生成请求并发派发出去。因此，在分析该代码的关键路径时，候选生成时长应按这段顺序相加，不能直接取最慢候选的一次耗时。把它重写成真正并发的版本可以成为工程优化，但应记录这是本书提出的改写，不应说论文原本已获得相应的时延收益。

其次，集成收到多个回答，并不等于它掌握了正确标签。固定公开实现中的 ScEnsemble 让模型返回候选编号，再返回对应候选的原文；Programmer 则另行生成代码并尝试执行\cite{aflow-code-2025}。选择已有答案与生成新检查是两种不同的责任。若几个候选沿用相同错误前提，聚合步骤可能只是在几种表达里选出最一致的错误。本例的 Programmer 又接收原题和集成分析，提供了一种可执行检查路径；它究竟调用了多少次模型、执行了几次代码、怎样处理错误，取决于算子的内部实现。顶层出现一次算子调用，只能说明有一次入口调用，不能替代底层模型请求清单。

最后，图中返回分支检验的是输出是否非空。原文 Appendix~A.4 展示的 Programmer 在连续执行失败后，返回的输出是字符串 \texttt{error}\cite{aflow-2410}。若按这两段展示代码直接组合，这个非空字符串也会进入采用分支。这个接口例子说明，执行状态与结果值需要分别传递，错误回执不能靠“是否有文字”识别。本书据此指出展示接口的组合风险，没有重建原实验中错误采用的发生次数。轨迹仍需保留检查状态、实际返回分支与最终任务验收。

\begin{principlebox}
\textbf{角色名说明设计意图，控制流说明实际行为。}
“评审者”可能只有一次无标签的模型判断；“并行求解”可能写成逐次等待；“验证通过”可能实际只是非空检查。分析一个工作流时，应从入口逐步列出等待关系、分支条件与算子内部调用，再给它贴机制标签。论文的流程图、示例代码与运行日志分别提供不同层级的证据，三者不能互相代填。
\end{principlebox}

论文 Figure~6 还展示了搜索如何在增加算子和修改提示之间交替前进。正文讨论的失败尝试包括一轮直接改写答案的评审，以及一轮过度强调题目中折扣信息的重述\cite{aflow-2410}。这些记录有助于理解搜索的具体动作，却不能把某一次分数下降全部归因于“评审有害”或“重述有害”。候选改变的是某个接口、某段提示和某条路径，验证集规模、采样波动及选取规则都影响观测差异。可靠的归因应回到保存的父子版本，并在同一组题上比较它们，而不是从搜索树上挑一条漂亮的上升曲线。

这也解释了为什么失败候选仍有保存价值。它可能暴露一个格式契约缺失，可能说明当前求解器不适合某种检查方式，也可能只是估计不够精确。若档案只保留最终赢家，后续搜索者无法分辨这些情况，只能重复尝试。若把失败候选写成“永远禁止此类结构”，又会把局部观测误升格成普遍规则。更可用的档案记录的是候选版本、改动描述、评价题目、逐题反馈和成本范围，随后由新的证据决定是否复用。停止规则同样需要这些记录：只看保留版本的最高分曲线，会把被拒候选隐藏起来。第\ref{sec:ch13-stop-output}节进一步区分停止搜索与选择交付版本；这不是AFlow原实验已经采用的附加规则。

\subsection{先画出值怎样流动，再给节点命名}
\label{sec:ch09-contractcase}

设想一个教学工作流：求解者读取编程题，生成候选程序；检查者运行公开样例；修订者根据错误修改程序；汇总者交付最终文件。这里没有本书新增的模型结果，只讨论这一设计能否被明确执行。若节点之间只传一段自然语言，“运行通过”可能指语法检查通过，也可能指公开样例通过；下游无法知道这一结论针对哪一个文件版本。即便四个角色各自按照提示认真工作，信息的歧义仍可能让整个流程出错。

因此，边上至少要传递三个不同对象：候选工件的版本、检查事件的状态，以及该事件实际支持的结论。程序版本应由内容或不可变标识确定，检查事件应记录命令、退出状态和输出，结论则注明覆盖的题目范围。公开样例全部通过，可以被表述为“该版本通过了这些样例”，不能直接压缩成“代码正确”。若下一节点又修改了文件，旧检查结果仍属于旧版本，不应自动附着到新版本上。

这个例子说明，类型契约不能只停留在“输入为字符串、输出为字符串”。较有用的契约还包含状态。检查者可能返回通过、失败、超时、工具不可用，或输出不能解析。将后三项全部编码为失败，会把运行环境问题混入算法质量；将它们编码为通过，则会放行没有证据的候选。执行层可以依据任务需要选择重试或后备路径，但证据层应保存原来的区别。下一节点收到的既是一个值，也是对这个值能否使用的限定。

错误还可能发生在边的另一端。假设修订者收到的是候选版本甲的失败日志，却在版本乙上继续修改。日志文字本身完全真实，仍然可能把修订方向带偏。只核对“是否有错误信息”不能发现这种问题；必须核对日志与候选之间的绑定关系。工作流图中的一条箭头因而不只是通信通道，也是对应关系：谁产生了这一观察，它描述哪个对象，接收者要用它作哪一个决定。

在设计接口时，可以把证据放在独立字段，把自然语言解释放在另一个字段。后者便于模型理解原因，却不能覆盖前者。若解释说“全部测试通过”，事件字段却记录超时，运行器应暴露冲突，而非取听起来更合理的一项。若解释无法由事件推出，可以保留为待验证假说。这样做并不会使模型推理自动正确，但能避免不同性质的信息在跨节点传递时被悄悄升级。

接口也需要注明允许修改的对象。一个评审节点可能只允许提出建议，另一个修订节点才有权写入文件；流程图中都画成一个方框，并不能说明权限相同。若评审者在检查前改了候选，原本计划测量的对象已经改变。若汇总者为了“润色”又重写最终代码，前面通过检查的工件可能根本没有被交付。最终返回应明确指向已验收的版本；需要另作修改时，就开启新的候选与检查记录。

这些要求使工作流从一张角色图变成可执行协议。角色名称仍然有帮助，它向读者解释分工；但真正决定可复核性的，是每条边携带的信息、版本和使用条件。AFlow 的非空返回条件就是一个小例子：只有读出这个条件，才能知道图中的检查者实际控制了哪一个决定\cite{aflow-2410}。没有这一步，再复杂的角色描述也不足以解释最终答案从哪里来。

\section{迁移表中一个容易漏掉的选择条件}
\label{sec:ch09-transferselection}

ADAS 的跨模型 ARC 实验并非随机抽取发现的全部结构。原文说明，它先根据 GPT-3.5 上的测试准确率选取表现最好的 \cited{3} 个候选，再把这些候选交给其他模型执行\cite{adas-2408}。这使读者能够考察选中结构在新模型上的表现，却也限定了结论的对象：结论针对这批按来源表现筛选的候选，不能直接推广为搜索产物普遍可迁移。来源测试集已经参与选取，因而也不再承担“完全未参与选择的最终来源评估”这一角色。

Table~3 的逐项比较比“迁移后一致超过基线”的概括更有信息。以 Structured Feedback and Ensemble Agent 为例，Claude-Haiku 上的点估计是 \cited{5.0}，同列 Self-Refine 是 \cited{6.3}；Claude-Sonnet 上分别为 \cited{38.7} 与 \cited{39.3}\cite{adas-2408}。这些点估计不能支持该候选在每个接收模型上均优于最强手工基线。其他候选在相应列可能更高，这支持“存在表现更好的候选”这一较窄陈述。这里也不能仅凭两条单独区间重叠与否宣布差异显著，需要与比较设计相符的差异估计。

迁移审计因而至少有两个不同的分母。一个是选中的候选数量，回答选中后有多少在接收方有效；另一个是搜索生成的全部候选数量，回答整个搜索过程发现可迁移资产的效率。只给前者，无法推断后者。若为了寻找迁移最强的结构又反复查看目标测试结果，目标测试集也会进入选择循环。应再保留一个未用于挑选的评价范围，或把结果明确标成探索性比较。

部署决策还要补上重新适配的费用。沿用来源结构、只调整接口、重新搜索，是三个不同的处理。沿用结构可能省去搜索但损失部分质量；重新搜索可能提高质量却耗时更长；调整接口则需要界定哪些改动仍算迁移，哪些已经改变结构。没有统一的定义，就会把一次重写包装成一次无损复用。第~\ref{ch08} 章对技能版本和依赖的要求，在这里同样适用于工作流版本与模型接口。

本书 E06 的原型尚未运行。本节是对公开论文和代码图的核验，不是本书的新实验。下一步若实现 E06，应先保留候选代码、算子实现和请求清单，再讨论结构是否值得搜索；否则即便得到了一个分数，也无法回答本章最初提出的预算问题。

\section{在质量、调用数与等待时间之间选择}
\label{sec:ch09-tradeoff}

三个量要分开报，因为它们不能互相换算：
\begin{itemize}
  \item \textbf{质量}：分数或准确率。单独报告它不能解释资源代价。
  \item \textbf{调用与token}：共同影响费用；单靠调用数不能确定账单，并行也可能增加等待。
  \item \textbf{等待时间}：决定能不能上线。串行评审不增加并行度，但拉长关键路径。
\end{itemize}

AFlow 的该成本图将测试 token 用量按模型价格折算\cite{aflow-2410}，
因此第二项只有一部分进了账，第三项没有进账，
而\emph{搜索本身}的成本没有包含在此测试费用比较中。
这些口径仍不足以独立证明``结构搜索比多采样划算''这一更窄结论——
E06 要补的就是这个对照。

% ---------- 5. 条件变化 ----------
\subsection{一次逻辑调用可能发生多少次实际请求}
\label{sec:ch09-occurrences}

对工作流计费时，计数对象应是发生过的调用事件，而不是图上画出的节点。设一个候选生成节点在循环里被调用多次，每次调用又可能因网络故障重试；图上仍然只有一个节点，服务端却可能收到多次请求。相反，一个节点可能命中本地缓存而不发出新请求。静态节点数只能描述结构规模，不能直接作为资源消耗。

可以为每次事件保存节点标识、所在迭代、父事件、逻辑请求编号和网络尝试编号。父事件帮助重建调用关系，逻辑请求编号把同一次决策的重试归在一起，网络尝试编号则保留实际派发次数。某次尝试没有收到响应，不代表它在服务端没有执行或没有费用。已知 token、请求次数、未知费用要分别保存；不能因为返回字段缺失就给它填零。

教学设计中，设首轮生成消耗一次逻辑模型调用，每轮评审一次，未通过时修订一次，最多允许两次修订。即使把生成、评审、修订各画成一个节点，最坏路径仍然是生成、评审、修订、评审、修订、评审，合计\calcnum{6}次逻辑调用。若每个逻辑调用最多允许\calcnum{3}次网络尝试，派发上界就是\calcnum{18}次。这里的数字来自规定的控制流，不是模型实测；若算子内部还有未计入的模型检查，上界就必须重新计算。

这个上界还有一个容易漏掉的前提：重试的层级已经统一。若外层允许三次尝试，而客户端内部每次也允许三次，最坏会变成每个逻辑调用九次派发。业务代码看见的“调用了一次客户端”并不等于服务端只收到一次请求。本书新版 E11、E12 的做法是关闭所用后端的内部重试，由外层记录每一次尝试；它解决的是当前实验的计账问题，不代表所有部署环境已经具有这一保证。

循环预算应与请求预算共同使用。只限制修订次数，无法限制一个修订节点内部再开一个无界循环；只限制模型调用，无法限制一次工具执行无限等待。实际运行需要分别说明模型请求、工具步骤、等待时间和输出规模的约束。某些约束可以在派发前检查，另一些只能在完成后得知。例如，不知道下一次真实输出长度时，事前的 token 估计不等于已实现的精确费用硬上限。

预算耗尽后的行为也属于结构定义。一个流程可以返回最后一个已检查版本，可以返回明确的未完成状态，也可以保存现场后终止。它不能在最后一刻绕过预算再请一个模型“总结一下”，然后把总结调用从账上删除。若最终答案来自尚未验收的版本，应清楚标记；如果直接把未完成序列按失败纳入成功率，应事先说明这是业务评价口径，而非运行器无意填出来的默认值。

这类规定看上去属于执行细节，却会改变搜索偏好。一个经常触发重试的结构，按成功返回计费时可能显得很便宜；一个擅长用超长上下文减少调用次数的结构，在只看调用数时可能显得最优。搜索器会利用评价口径中的缺口，无须具有欺骗意图。把事件账写完整，才能让“更好的结构”对应到真正要优化的资源条件。

在各支路均须完成、聚合随后开始，且节点耗时按同一种执行资源计量的条件下，不能直接把节点耗时之和当成并行段等待时间。设并行段有 $m$ 条支路，
第 $j$ 支的节点耗时为 $t_{j,i}$，聚合耗时为 $t_{\mathrm{agg}}$，
并发上限为 $c$，则该段的等待时间不可能短于

\begin{equation}
T_{\mathrm{par}} \;\geq\;
\max\Bigl(\;\max_{j}\sum_{i} t_{j,i},\;\; \frac{1}{c}\sum_{j}\sum_{i} t_{j,i}\Bigr)
\;+\; t_{\mathrm{agg}},
\label{eq:ch09-latency}
\end{equation}

这是上述资源与依赖假设下的下界；缓存、提前退出或异构资源改变工作量时，须重新建立模型。
\textbf{并发充分时最长支路提供约束，串行时总工作量提供约束；一般并发条件下，下界不保证可达}。
在单资源串行、无额外等待的理想条件下，工作时间确实退化为各支路节点耗时之和；并发充分时则由最长支路约束。
把并行支路的耗时相加后直接称作用户等待时间，
混淆的是累计执行耗时与端到端墙钟；它们也都不是调用次数。

\subsection{把并行下界与可以达到的时延分开}
\label{sec:ch09-scheduling}

式~\eqref{eq:ch09-latency} 给出的下界有两个来源：一条有依赖的支路必须等前一步结束，总工作量也必须被有限资源执行完。这两个条件都成立，不表示取二者较大值就一定能安排出那样快的调度。任务是否可切分、服务是否同时可用、某些节点是否争用同一个资源，都会使实际时间更长。

考虑三个互不依赖、不可中断的教学调用，每个都占用同类资源\calcnum{2}个时间单位，并发上限为\calcnum{2}。最长支路为\calcnum{2}，总量除以并发为\calcnum{3}，所以公式下界为\calcnum{3}。实际只能先执行两个调用，第三个等其中一个结束后开始，最早在\calcnum{4}结束。聚合若再占用一个时间单位，最终为\calcnum{5}，而不是由下界误算出的\calcnum{4}。这不是公式失效，而是把必要条件当成了可达方案。

若调用时长不同，派发顺序也可能改变结果。设三个不可中断调用耗时依次为\calcnum{3}、\calcnum{2}、\calcnum{2}，资源仍为两份。先派发两个短调用，再派发长调用，全部完成需\calcnum{5}；先派发长调用和一个短调用，再补另一短调用，只需\calcnum{4}。这组数字是确定性排程演示，不预测真实 API 时延。真实服务中时长事前往往未知，因此“先运行最长任务”本身还需要估计，不能把事后最优排程当作在线算法已经做到的事。

提前返回会进一步改变比较对象。若策略是第一个答案返回后立即交付，它测的是最早完成者的质量；若等全部候选再投票，它测的是聚合后质量。两者的等待时间与选择规则同时变化。取消尚未完成的请求是否成功、取消前是否已经产生费用，也需要单独记录。不能把最终只采用一个答案写成只消耗一次生成，更不能把未采用的支路从调用账中抹掉。

并发还会通过环境反过来影响每个节点的耗时。离线单独测得的调用时间，并不必然等于多个调用竞争服务时的耗时。限流、重试、队列和共享带宽都可能改变分布。因此，可以用串行测量建立粗略预测，但最终应同时保存派发时刻、首个响应时刻和完成时刻，以区分排队与生成。只有最终总耗时，通常无法判断慢在结构依赖还是服务拥塞。

部署关心的也未必只是平均时间。一个平均较快但偶尔等待很久的结构，可能违反响应期限；一个稍慢却稳定的结构，可能更适合交互任务。报告平均值、分位数或超时比例时，应使用一致的任务范围和超时处理规则。把超时任务从耗时样本里删除，会使最不稳定的结构看起来很快；全部记为超时上限又是另一种截断口径，应注明而非假装观测到了真实结束时刻。

工作流搜索因此可以有不同的优化目标：固定响应期限下最大化有效完成率，固定费用下最大化质量，或在几个指标之间保留可选方案。没有必要把它们强行压成一个“综合分”。如果确实使用加权目标，就应说明权重由谁决定、适用于何种业务范围。优化器只是在追逐给定的评价函数，不能替读者决定多等一分钟与多答对一道题之间的价值关系。

\section{best-of-N 的上限并不是一个免费的选择器}
\label{sec:ch09-selector}

在讨论扁平采样对照时，必须先问“best”由谁判断。若使用任务的隐藏答案挑出正确候选，就得到一个有真值访问权的上限；部署系统通常没有这种权限。若让另一个模型选择，就引入新的调用、偏差和错误；若使用可执行测试，则选择器只能依据实际可见的测试，而非未公开的最终验收集。三种方法可以分别有研究价值，但不应共用一个名称后直接比较。

设每个候选独立正确的概率为 $p$，共生成 $n$ 个候选。至少有一个正确候选的概率是
\begin{equation}
 P(\text{存在正确候选})=1-(1-p)^n.
 \label{eq:ch09-oracle}
\end{equation}
这是独立假设下的生成覆盖率，不是最终交付的正确率。完美选择器才能在所有存在正确候选的情形下都选对。若选择器在“至少有一个正确候选”的条件下选对的概率为 $q$，且它只能选择而不能另写一个答案，则最终正确率为 $q[1-(1-p)^n]$。这里 $q$ 的定义已经包含不同候选组合的难度，无须另外假定选择错误与生成错误独立。

教学设定 $p=\calcnum{0.6}$、$n=\calcnum{3}$ 时，独立覆盖率为\calcnum{0.936}；若条件选择率为\calcnum{0.7}，交付正确率只有\calcnum{0.6552}。这比单次生成高多少，需要连同三次生成与选择的成本一起判断。若三个候选的正确事件完全相同，覆盖率仍是\calcnum{0.6}，不能使用独立公式。这里比较的是两种明确的概率模型，不是断言所有同模型候选都完全相关。

选择器还可能执行改写。那时它有机会在所有候选都错误时修正答案，也可能把一个原本正确的答案改坏。最终概率需同时计入“有正确候选时成功”和“没有正确候选时成功”两个条件，前述简单乘积不再完整。由此可以看出，Ensemble、投票、选择、综合重写并不是同一个操作；报告“采用集成”而不说明内部规则，很难确定比较的到底是什么。

对于代码任务，公开测试提供了可执行选择信号，却不是万能裁判。多个错误程序可能都通过有限样例；更长的候选可能通过样例却引入新的边界错误。若同一个公开测试集被反复用于修订，它已参与搜索。最终评价应保留一个用途明确的验收范围，或坦率承认比较仅反映公开测试上的表现。选择器读取隐藏验收结果的次数，也需要进入查询预算。

一个可执行的等预算基线可以这样定义：先规定总请求或 token 上限，再在这个上限内分配生成与选择预算。选择器不能在生成用完预算后免费出现。若候选较长导致最后一次选择无法容纳全部文本，应按预定规则截断、分批或终止，并记下实际输入范围。结构方法若有类似聚合开销，也采用相同原则。这种比较未必最强，却比把真实方法与免费真值选择器混为一谈更清楚。

真值上限仍然值得报告。它可以区分“没有生成正确候选”和“生成了但没有选中”两类瓶颈。若上限很高而实际选择很差，增加更多生成未必是优先方向；若上限也低，改进选择器能带来的空间就有限。不过这属于诊断分解，不能把有真值帮助的上限当成可部署对手，然后据此宣布某个结构已被现实基线替代。

\section{同一张费用表可以回答两个不同的问题}
\label{sec:ch09-costcomparison}

回到 AFlow Appendix~D。该表的具体任务是 HumanEval，表中 DeepSeek 执行由 GPT-4o-mini 反馈搜索出的工作流，分数为\cited{0.9390}，执行费用为\cited{0.0291}美元；GPT-4o 直接调用分数为\cited{0.9389}，费用为\cited{0.6371}美元\cite{aflow-2410}。按展示的四位小数重算，费用比约为\calcnum{4.57\%}，与论文正文报告的\cited{4.55\%}存在数值不一致，差异来源尚未确定。不能把重算后的最后一位当成新测量，也不能据两个点值接近就证明统计等效。

如果读者关心“较便宜模型配合结构能否接近较贵模型”，这是一组有关的比较。但它同时改变了模型与结构，不能单独隔离结构的贡献。看同表 DeepSeek 自己的直接调用，分数是\cited{0.8860}，费用是\cited{0.0127}美元\cite{aflow-2410}。与它相比，前述工作流的点估计增加\calcnum{5.30}个百分点，执行费用约为\calcnum{2.29}倍。换一个合理的参照，结论从“较低费用接近强模型”变成“同一模型上更多执行费用换来更高点估计”；两句话可以同时成立。

这也提醒读者不要把展示口径混入公式。原表分数栏虽然带百分号，单元格使用的是零到一的数值。本书在相减后明确乘以百分之百，并用百分点描述差值。费用列属于该表报告的测试执行范围，不应擅自改称每次 API 请求的单价或每题费用。只有知道分母与汇总范围，才能把它换算到另一个使用场景。模型今日价格与论文当时价格也不应混用，本章不将历史价格当作当前报价。

生命周期判断还需要搜索费。设方法甲需要一次搜索投入 $S$，每个可比部署单位花费 $a$；基线乙不需要该搜索投入，每单位花费 $b$。在同等质量要求、相同覆盖范围且 $b>a$ 的教学模型下，甲在 $K$ 次复用后费用更低的条件是
\begin{equation}
 S+Ka<Kb,\qquad K>\frac{S}{b-a}.
 \label{eq:ch09-amortization}
\end{equation}
这里 $S$ 应包括搜索失败候选和用于选择的验证，不只是最终赢家的生成。若有持续维护或重新适配，需再加相应项；若甲每次都更贵，仅靠复用次数增加并不会在这个纯费用模型中反超。

以明确的教学单位举例，若搜索费为\calcnum{120}，每次使用甲花\calcnum{2}、乙花\calcnum{5}，则\calcnum{40}次时总费相等，严格低于要从整数复用量\calcnum{41}起，与式~\ref{eq:ch08-integer-threshold}的取整区别一致。这个例子与 AFlow 的实测无关，因为本章没有它所需的完整搜索费与相同部署单位。若部署只运行十次，是否仍采用甲可能取决于质量提升，但那已是带质量约束的决策，不能继续声称单靠省钱回本。

摊销还受资产寿命限制。模型接口改变、任务分布变化或依赖升级，都可能使旧工作流需要重新验证。若预期复用量很大而实际有效寿命很短，按无限复用作出的经济判断会失效。应记录何时重新验证、何时退役，以及退役前实际服务了多少任务。第~\ref{ch08}章的技能生命周期在结构层继续成立，只是需要保存的资产从函数扩展为代码、提示、路由和执行环境的组合。

由此，费用报告最好允许读者重建不同问题的答案：同一后端的结构增量、不同后端的部署替代，以及包含发现过程的生命周期成本。它们共享部分记录，却不共享全部分母。没有完整数据时，可以保留已经可回答的较窄问题，明确另一项未知；没有必要为了得到一个整齐的胜负结论，把缺失费用默认为零。

\section{把 E06 设计成能够失败的比较}
\label{sec:ch09-e06design}

E06 仍未运行，下面是设计要求，不是新的实验结果。首先应限定主张：要研究的是结构本身相对等预算采样的增量，还是搜索流程相对人工设计的整体价值。前者需要尽量对齐后端、工具权限和实际资源；后者还要记录人工或自动设计所使用的信息与投入。两个问题可以连续研究，但不能用一次不完整对照同时回答。

一个最小比较可以保留直接调用、固定扁平采样和搜索所得结构三种处理。直接调用说明基础能力，扁平采样把额外推理资源投入候选生成与选择，结构处理允许有条件的信息流。三者使用相同的最终任务与验收规则。搜索阶段只使用允许的训练与验证范围，最终结构冻结后再进入评价。若扁平基线也需要调采样数或选择提示，这些调参同样属于发现成本。

“等预算”必须有主约束。固定调用数不保证 token 相同，固定 token 不保证响应时间相同，固定历史美元也不保证当前服务成本相同。可以选择与实际用途最接近的一项作为主要上限，同时把其余指标完整报告。在预算上限相同而实际消耗不同的情况下，应同时报上限与实耗，不能仅凭上限相同就声称每个方法花了完全一样的钱。

运行顺序应尽量避免服务条件与处理绑定。总是先跑基线、最后跑结构，可能把时段限流或后端更新混进差异。可以使用预先固定的交错或随机顺序，并记录后端标识、时间与异常。若预算只允许很小规模，仍可先检查装置是否正确，但应把它标成工程试跑；小试暴露的格式错误可以修复，却不能挑出漂亮题目就直接升级成确认性结果。

分析时先看共同覆盖的任务和未完成原因。某结构在困难任务上经常耗尽预算，如果只报告成功完成的题，会得到过于乐观的质量与时延。业务完成率可以把超预算视为未完成，同时另报已完成题的质量；缺失日志或服务失败则按预定方案处理，并保留敏感性分析空间。所有这些分母都应在看见结果前明确，否则“同等预算”会在筛选数据时失去意义。

若结果显示扁平采样在设定范围内能够达到相近质量与资源表现，这会削弱采用该复杂结构的理由，却不能证明所有结构搜索无效。若搜索结构更好，也需要检查优势是否来自更多验证查询、不同信息输入或未计入的选择器。一个允许自身失败的实验，应当事先写明哪些观察会使当前主张收窄，而不是任何结果都能被解释成方法成功。

最后应交付能够重建比较的工件：冻结结构与算子实现、逐请求账、逐题结果、选择规则、预算终止记录，以及搜索阶段的候选档案。漂亮的流程图帮助理解设计，完整工件让别人判断结论是否成立。对本书而言，在 E06 尚无这些实跑产物时，保持未运行状态比填上一组借来的数字更有价值；公开论文提供机制线索，本书自己的判断仍应由相应的比较承担。


% ---------- 6. 选择原则 ----------
\subsection{条件变化：换域、换模型与换预算}
\label{sec:ch09-conditions}

\begin{itemize}
  \item \textbf{换域。}ADAS 从 MGSM 数学域发现的候选迁移到 GSM8K、GSM-Hard 数学任务时，相对所列手工基线最佳值的差为 \cited{25.9} 与 \cited{13.2} 个百分点；数学迁移到 DROP/MMLU 是另外两列，不能把前两个数字归到阅读理解。原文独立域内搜索的 DROP 与 MGSM 增益又是不同实验\cite{adas-2408}。
  \item \textbf{换模型。}AFlow Table 2 为每个执行模型列出 IO、CoT 等基线，同时比较两种来源执行模型搜出的工作流；换后超过基线与不如目标模型专门搜索的版本可以同时成立。ADAS Table 3 也列有各目标模型的手工基线，不能声称两者都缺基线。原报告支持具体范围内的迁移比较，不构成相互矛盾的普遍规律\cite{aflow-2410,adas-2408}。
  \item \textbf{换预算。}本章仍需单独设计``同等预算下的扁平采样''对照\cite{aflow-2410}。
  \item \textbf{通信截断。}多 Agent 结构应明确每条通信边传递哪些内容。
        一旦通道有长度上限，或被一次摘要压缩过，
        被丢掉的那部分就是\emph{结构自己引入的信息损失}：
        其影响还与节点如何利用信息有关，需要与结构干预一起解释。
        因此应保存每个节点实际接收的内容与长度，
        仅有字符数还不足以区分结构、信息来源与通道容量的作用。
  \item \textbf{评审者相关错误。}加一个评审节点既能拦下真错误，
        也能把正确答案否掉，两种错误的代价并不对称。
        更要紧的是：评审者与求解者共享模型、数据或提示时，错误可能相关，不能仅由角色名称假定独立；
        若使用基于独立正确事件的投票概率模型，就需检验相应假设；考虑相关性的模型并不要求独立。
        若错误正相关却按独立处理，投票稳定性可能被高估——
        这与第 \ref{ch04} 章讨论的``自评不是独立反馈''是同一个问题，
        只是在那里发生在反馈层，在这里发生在结构层。
  \item \textbf{并发受限。}并行的路数受到并发上限约束时，
        等待时间还取决于排队和派发顺序。并发上限降到一时，各支路只能依次执行；
        在调用集合和返回规则相同的条件下，请求数可以不变，等待时间却接近各支路之和。
        \textbf{即使分数、调用数与逻辑结构相同，也不能据此推断执行资源条件相同}——
        这正是``只报分数''最彻底的失效方式。
  \item \textbf{任务变短。}结构搜索是在一批任务上做的，
        而任务的典型长度会变。任务变短之后，
        固定开销（分发、评审、聚合）占的比例上升，
        为长任务搜出的结构在短任务上可能净亏。
        因此结构必须与它针对的任务长度一起报告，
        不能只报``在某个基准上搜出来的''。
\end{itemize}

\begin{principlebox}
\textbf{迁移收益与专门适配收益是两个问题。}
迁移后的工作流可以超过目标模型的直接调用，同时仍落后于为目标模型重新搜索的工作流。这两条事实不矛盾。目标模型基线用于计算迁移增量，目标模型专门搜索用于比较重新适配的价值；不同主张需要不同对照。E03c 已完成本书代码经验的换模型比较，但不能替代 AFlow 或 ADAS 的工作流迁移实验。
\end{principlebox}

\section{选择原则}
\label{sec:ch09-choice}

\begin{enumerate}
  \item \textbf{报分数时必须同时报调用数与等待时间。}
        三者不可换算；只报第一个等于无法判断这个结构值不值。
  \item \textbf{搜索成本要和使用成本进同一张账。}
        “费用为基线的 \cited{4.55\%}”是测试执行费用比，不是“便宜4.55\%”；若讨论生命周期，应另计搜索及复用范围。
  \item \textbf{加角色之前先问它替代了什么。}
        增加评审节点带来纠错机会，也带来调用数与等待时间；
        如果它只是把一次错误变成两次调用，那就是净亏。
  \item \textbf{把迁移增量与重新适配分别比较。}
        换后模型的直接调用基线，以及用它自己搜一遍的基线。
        前者支持迁移增量，后者支持迁移与重新搜索的比较，不应把两者机械地设为所有迁移主张的必要条件。
  \item \textbf{根据主张设置单次调用与等预算采样对照。}
        单次调用可以作为基础参照，但不足以隔离更多计算的影响；结构搜索也没有几乎必胜的保证。
  \item \textbf{用消融解释整体比较，并检查部件交互。}
        在保留必要接口的条件下删除节点或边，观察其条件贡献；有互补机制时还需联合消融，不能由逐项删除分配唯一功劳。
        \emph{若消融差异没有得到支持，应继续检查精度与触发覆盖，不能直接认定该部件毫无价值}——
        它们既可能是冗余路径，也可能只是尚未被这批题触发。
        整体比较可以支持整套流程的增量，部件归因则需要相应消融，二者的证据范围要分别说明。
  \item \textbf{只有当结构带来可复核的增量时才保留复杂度。}
        应预先说明希望改善的质量或资源量、可容忍的其他损失与实际重要阈值，再检查证据是否支持这一取舍；任何方向的显著变化都不能自动算成收益。
        证据不足时可以暂缓采用复杂结构，也应保留不确定性，不将一次未显著消融等同于无效。
\end{enumerate}

% ---------- 7. 反例 ----------
\begin{pitfallbox}
\textbf{反例一：把结构变复杂当成找到了更好的结构。}
图~\ref{fig:ch09-workflow} 里三种组织可以拿到同一个分数。
复杂度上升与分数上升同时发生时，要能说明是哪一个节点起了作用——
否则增加的可能是纯粹的冗余路径。

\textbf{反例二：把搜索成本排除在比较之外。}
\cited{4.55\%} 的成本比只统计测试时推理 token，
最多 \cited{20} 轮搜索、每轮 \cited{5} 次验证所产生的搜索开销不在其中\cite{aflow-2410}。
该费用比不含搜索，加入搜索之后需要实际投入与部署复用量才能计算，不以当前测试费用表代填。

\textbf{反例三：把跨模型能用当成跨模型无损。}
AFlow 与 ADAS 都提供了目标模型上的基线比较。迁移超过基础参照，并不保证优于目标模型专门搜索，也不保证每个候选在每个模型上都有正增量。应逐项读取表格，而不是把两篇论文概括为互相矛盾；本书 E03c 不负责裁决外部工作流实验。

\textbf{反例四：把``比单次调用强''当成``比等预算采样强''。}
基线里已有 \cited{5} 答案自一致性与 \cited{3} 答案 \cited{5} 票这类多尝试方法\cite{aflow-2410}，
但缺``同等预算下 best-of-N''这一组\cite{aflow-2410}。
E06 的证伪条件写的就是这一组——它还没被跑过。
\end{pitfallbox}

\section*{习题}
\begin{enumerate}[label=9-\arabic*.,ref=9-\arabic*]
  \item\label{ex:ch09-1} 取图~\ref{fig:ch09-workflow} 的中间那条（带评审），
        写出它在一次任务上的调用数、等待时间与失败路径。
        若教学假设评审有30\%的误放概率，还需哪些错误率、费用与基线信息才能计算净收益？
  \item\label{ex:ch09-2} AFlow 的算子集合里有 Ensemble。若把并行路数从3路提到9路（教学设定），
        调用数、等待时间与错误相关性分别受哪些条件影响？是否有一个量必然不变？
  \item\label{ex:ch09-3} ADAS 的 archive 服务于搜索，和部署经验用途不同。
        先列出三类必需的来源与评价信息，再说明恢复候选和跨搜索采用还依赖哪些条件。
  \item\label{ex:ch09-4} 设计一个最小实验，使得``同等预算下扁平采样能否复现结构搜索的收益''可判。
        说明你会固定哪个预算量（调用数、token、墙钟还是美元），以及为什么。
\end{enumerate}

\section*{常见误解}
\begin{itemize}
  \item \textbf{``多 Agent 就是多个模型实例''。}
        在工作流表示里，角色是节点、通信是边；
        角色可以由同一个模型的不同提示承担。本章不用实例数定义多 Agent。
  \item \textbf{``搜出来的结构比手设计的强，所以搜索有用''。}
        这句话漏掉的是预算。单次调用是有用基础参照，等预算多采样进一步检查收益能否仅由计算量解释。
  \item \textbf{``并行不增加等待时间，所以并行是免费的''。}
        等待时间取决于支路时长、并发限制、排队与聚合，调用与token也需要单独结算。图~\ref{fig:ch09-workflow} 右栏画的就是这个。
  \item \textbf{``换模型还能跑，说明结构是可移植的''。}
        AFlow Table~2 同时显示部分迁移版本高于目标模型基础参照、低于专门搜索版本\cite{aflow-2410}。
        可运行、有正迁移增量、达到专门适配水平是三个不同要求，后两者不能互相冒充。
\end{itemize}

\section*{进一步阅读}
\begin{itemize}
  \item \cite{aflow-2410,aflow-code-2025}：原文目标函数（本章式~\eqref{eq:ch09-aflow}）与MCTS四步；Table~1对应表~\ref{tab:ch09-aflow}，Appendix~D给出\cited{4.55\%}的测试费用比。固定代码补充选择参数、分数变换与返回行为；该测试费用不含搜索投入。
  \item \cite{adas-2408}：代码定义agent与archive，以及摘要、§4的跨域迁移。来源候选按测试成绩筛选；任务或基线不同，不能直接相减当作纯迁移损失。
  \item 第5章：两种拟合目标与E03c的换模型比较及未决范围。经验迁移不替代工作流迁移的证据。
  \item 第7章：提示与上下文优化，对应表~\ref{tab:ch09-space}左栏。
  \item 附录B：当前实验台尚未实现E06；本章保留设计与外部来源，不将未实现写成永久限制。
\end{itemize}
